\section{Components and the pooled reference}
\label{sec:components}

\paragraph{Ablations of V0 (the plan's claim D1) and of V2 (exploratory).}
The plan's ablations removed components from V0, because they were
written before V2 existed. They were run on held-out accounts in a second,
separately logged run, after the main results were known
(Table~\ref{tab:ablations}, Appendix~\ref{app:tables}). The V2 ablations
were added later, on both splits (development results in
\texttt{paper/results/exploratory/s6-v2-ablations/}). Below: each variant's
skill against its full model on held-out accounts, account mean with the
pooled skill in brackets (positive: removing the component lowers CRPS).

\begin{center}\scriptsize
\begin{tabular}{l ll ll}
\toprule
 & \multicolumn{2}{c}{in V0} & \multicolumn{2}{c}{in V2}\\
\cmidrule(lr){2-3}\cmidrule(lr){4-5}
Component removed & \hB\ d & \hC\ d & \hB\ d & \hC\ d\\
\midrule
adaptive noise & \R{ho/ablp/all/abl-fixed-noise/30/skill@3} {\tiny(\R{ho/ablp/all/abl-fixed-noise/30/pooled_skill@3})} & \R{ho/ablp/all/abl-fixed-noise/90/skill@3} {\tiny(\R{ho/ablp/all/abl-fixed-noise/90/pooled_skill@3})} & \R{xh/s6/V2-noadapt-vs-V2/30/skill_mean@3} {\tiny(\R{xh/s6/V2-noadapt-vs-V2/30/skill_pooled@3})} & \R{xh/s6/V2-noadapt-vs-V2/90/skill_mean@3} {\tiny(\R{xh/s6/V2-noadapt-vs-V2/90/skill_pooled@3})}\\
regime detector & \R{ho/ablp/all/abl-no-detector/30/skill@3} {\tiny(\R{ho/ablp/all/abl-no-detector/30/pooled_skill@3})} & \R{ho/ablp/all/abl-no-detector/90/skill@3} {\tiny(\R{ho/ablp/all/abl-no-detector/90/pooled_skill@3})} & \R{xh/s6/V2-nodetector-vs-V2/30/skill_mean@3} {\tiny(\R{xh/s6/V2-nodetector-vs-V2/30/skill_pooled@3})} & \R{xh/s6/V2-nodetector-vs-V2/90/skill_mean@3} {\tiny(\R{xh/s6/V2-nodetector-vs-V2/90/skill_pooled@3})}\\
spend floor & \R{ho/ablp/all/abl-no-floor/30/skill@3} {\tiny(\R{ho/ablp/all/abl-no-floor/30/pooled_skill@3})} & \R{ho/ablp/all/abl-no-floor/90/skill@3} {\tiny(\R{ho/ablp/all/abl-no-floor/90/pooled_skill@3})} & \R{xh/s6/V2-nofloor-vs-V2/30/skill_mean@3} {\tiny(\R{xh/s6/V2-nofloor-vs-V2/30/skill_pooled@3})} & \R{xh/s6/V2-nofloor-vs-V2/90/skill_mean@3} {\tiny(\R{xh/s6/V2-nofloor-vs-V2/90/skill_pooled@3})}\\
lumps / sparse spending & \R{ho/ablp/all/abl-no-lumps/30/skill@3} {\tiny(\R{ho/ablp/all/abl-no-lumps/30/pooled_skill@3})} & \R{ho/ablp/all/abl-no-lumps/90/skill@3} {\tiny(\R{ho/ablp/all/abl-no-lumps/90/pooled_skill@3})} & \R{xh/s6/V0-vs-V2/30/skill_mean@3} {\tiny(\R{xh/s6/V0-vs-V2/30/skill_pooled@3})} & \R{xh/s6/V0-vs-V2/90/skill_mean@3} {\tiny(\R{xh/s6/V0-vs-V2/90/skill_pooled@3})}\\
the filter & -- & -- & \R{xh/s6/cp-nokf-vs-v2/30/skill_mean@3} {\tiny(\R{xh/s6/cp-nokf-vs-v2/30/skill_pooled@3})} & \R{xh/s6/cp-nokf-vs-v2/90/skill_mean@3} {\tiny(\R{xh/s6/cp-nokf-vs-v2/90/skill_pooled@3})}\\
\bottomrule
\end{tabular}
\end{center}

\noindent Held-out origins: \R{ho/abldata/eligible_origins} per horizon for
V0's ablations and \R{xh/s6/V2-nodetector-vs-V2/30/forecasts} for V2's.

\begin{itemize}
\item \textbf{Removing the detector hurts both models, V2 far more.} The
  sign is the same as for V0; the magnitude is not. A diagnostic on
  development accounts (exploratory; Appendix~\ref{app:detector}) traces
  it. The whole V2 effect sits in accounts with detected income and an
  income alarm; elsewhere removing the detector leaves V2's forecasts
  unchanged. In Berka, the first income-tagged record is typically a
  small fraction of the later pay (median \R{xd/s9/mech/pay1-over-last/median}\%),
  and the income alarm, a median \R{xd/s9/mech/first-alarm-age/median} days
  into the history, is what moves the expected paycheck $m$ to the pay
  level. Without it, $\hat m$ stays at a median
  \R{xd/s9/mech/m-nodet-over-last/median}\% of the last paycheck, and V2's
  cone, whose spending model ignores the spend rate, projects most of the
  pay as missing. The V2 ablation therefore measures the detector's role
  in initializing the income state, not its response to later regime
  changes.
\item \textbf{An engine flaw, fixed after the freeze.} The cause is the
  income channel's outlier gate. After a small first income record, the
  real paychecks look like outliers and skip the paycheck update; the
  income noise estimate then inflates until they pass the gate, by which
  time the gain on $m$ is too small for it to catch up (median gain
  \R{xd/s13/trace/before-fix/income-alarm/nodet/median_k_last12} without the
  detector against \R{xd/s13/trace/before-fix/income-alarm/det/median_k_last12}
  with it; Appendix~\ref{app:detector}). We added a gate escape: with the
  income detector off, two consecutive income observations gated on the
  same side start a new pay level, as an alarm would. With the detector on,
  the engine is unchanged: the frozen V0 and V2, rerun with the fix on the same
  \R{d/pf-accounts} development accounts as the earlier post-freeze check
  (Appendix~\ref{app:identity}), differ in
  \R{xd/s13/fix/V2/90/forecasts_changed} of \R{xd/s13/fix/V2/90/forecasts}
  forecasts per horizon, and the frozen results stand. With the fix,
  removing the detector costs V2 an account-mean skill of
  \R{xd/s9p/after-fix/all/V2-nodet-vs-V2/90/skill_mean} at \hC\ days on
  development (pooled \R{xd/s9p/after-fix/all/V2-nodet-vs-V2/90/skill_pooled}),
  instead of \R{xd/s6/V2-nodetector-vs-V2/90/skill_mean}.
\item \textbf{The spend floor has no effect in V2}, whose spending model
  does not use the rate the floor acts on.
\item \textbf{Fixed noise priors are marginally better} in both models.
\item \textbf{Removing V2's sparse spending} (V0) is much worse.
\item \textbf{Removing the filter} (CP-noKF) is slightly worse.
\end{itemize}
V0's lumps row and V2's sparse-spending row are different components that
play the same role, the heavy tail of spending.

\paragraph{B-GLOBAL, the pooled-data reference (exploratory).} B-GLOBAL
is trained on \R{dev/data/accounts_run} other accounts, data a local-first
application cannot collect, though it could run a model trained
elsewhere. It has the lowest CRPS of any method on held-out data
(Table~\ref{tab:pairwise}). Its training used development origins of all
dates, a calendar look-ahead for early held-out origins. Retrained without
it, quarter by quarter on earlier development data only (G-cf), it loses
little: pooled skill against the original
\R{xh/s7-cf/B-GLOBAL-cf-vs-B-GLOBAL/30/skill_pooled} and
\R{xh/s7-cf/B-GLOBAL-cf-vs-B-GLOBAL/90/skill_pooled} at \hB\ and \hC\ days
(account mean \R{xh/s7-cf/B-GLOBAL-cf-vs-B-GLOBAL/30/skill_mean} and
\R{xh/s7-cf/B-GLOBAL-cf-vs-B-GLOBAL/90/skill_mean}). It still has lower
CRPS than B-EQ: pooled skill
\R{xh/s7-cf/B-GLOBAL-cf-vs-B-EQ/30/skill_pooled} and
\R{xh/s7-cf/B-GLOBAL-cf-vs-B-EQ/90/skill_pooled} (account mean
\R{xh/s7-cf/B-GLOBAL-cf-vs-B-EQ/30/skill_mean} and
\R{xh/s7-cf/B-GLOBAL-cf-vs-B-EQ/90/skill_mean}). Its advantage survives the
removal of calendar look-ahead. Appendix~\ref{app:bglobal} gives the
details and its weaker calibration.
